% =====================================================================
% Experimental design.
% Adapted from RESEARCH_PLAN.md §5 (models), §6 (axes), §8 (grid, ground
% rules, datasets).
% =====================================================================
\section{Experimental Design}
\label{sec:design}

\subsection{Models}
\label{sec:design-models}

Model selection follows three criteria: a shared architectural family
(decoder-only, RoPE, GQA) so the mask logic transfers unchanged; two
independently-trained families so parity claims cannot rest on a
single-lab artifact; and public reproducibility for every primary claim.

\begin{table}[t]
\centering
\caption{Models. Both base and instruct variants of each primary model are
    evaluated; the base/instruct pairing supports the mechanism probe in
    Section~\ref{sec:results-quality}. Ablation and scale-trend rows are
    listed here for reference but are not populated in the current
    results (see the note below the table).}
\label{tab:models}
\begin{tabular}{@{}llll@{}}
\toprule
Model & Role & Reproducible & Notes \\
\midrule
Qwen2.5-7B (base, instruct)     & Primary   & yes & Comparable to prior hypothesis runs \\
Llama-3.1-8B (base, instruct)   & Primary   & yes & Independent training family \\
Qwen2.5-1.5B                    & Ablation  & yes & Full mask $\times N \times \ell$ grid \\
Qwen2.5-14B (or 32B)            & Scale     & yes & Scale trend on affordability \\
YandexGPT-5-Lite-8B             & Appendix  & no  & Third-family confirmation \\
\bottomrule
\end{tabular}
\end{table}

\todo{The primary-model quality results reported in
Section~\ref{sec:results-quality} land only on Qwen2.5-7B (base and
instruct); the Llama-3.1-8B rows and the 1.5B/14B ablation-and-scale
sweep were scoped but not run. Any claim about Llama, sub-7B, or
$>$7B behaviour is deferred to a future evaluation cycle.}

All runs use bf16 precision, greedy decoding for quality metrics, and
seed control committed to provenance. Batch size is 1 for the current
harness; lifting this constraint via FlashAttention varlen is a
future-work item (Section~\ref{sec:limitations}).

\subsection{Datasets}
\label{sec:design-datasets}

Datasets are chosen by function, not by fashion. A parity-detection
benchmark and a failure-envelope benchmark are different instruments and
must not be mixed.

\begin{table}[t]
\centering
\caption{Datasets. Every parity-detection dataset has its seed variance
    measured under \texttt{full} and is admitted only if the variance is
    smaller than $\Delta$. Datasets marked ``no-doc'' in the notes column
    carry no retrieved passages and therefore collapse every mask variant
    to plain causal attention; they are retained only as harness
    smoke-tests.}
\label{tab:datasets}
\begin{tabular}{@{}p{0.28\linewidth}p{0.36\linewidth}p{0.28\linewidth}@{}}
\toprule
Function & Dataset & Notes \\
\midrule
Parity detection (single-hop)    & KILT-NQ, \texttt{synthetic\_qa} & With retrieved passages \\
Parity detection (no-doc)        & NQ-Open                         & No passages; smoke-test only \\
Multi-hop (natural)              & HotpotQA, MuSiQue               & Failure-envelope side \\
Cross-document stress (synthetic)& \texttt{composition}, \texttt{ddi\_stress} (fact\_chain, fact\_aggregation, contradiction, coreference, needle) & Family-controlled multi-hop probe \\
Comparability with prior work    & $\geq 1$ dataset PCW / FiD report on & Overlap with parity row \\
\bottomrule
\end{tabular}
\end{table}

\todo{TriviaQA, MoNaCo, Open RAG Bench, and the LitM-style multi-doc QA
protocol were listed as target loaders in the plan but not implemented in
time for the results reported here; when they land they extend
Section~\ref{sec:results-quality} and Section~\ref{sec:failure-envelope}
without changing this section's function-split framing.}

\subsection{Design axes}
\label{sec:design-axes}

The full study is organized along six axes; the sweep design in
Section~\ref{sec:results-quality}--\ref{sec:failure-envelope} does not
enumerate their cross product.
\begin{description}
    \item[Axis M --- mask.] \texttt{full}, \texttt{ddi\_qf},
          \texttt{ddi\_strict}, \texttt{ddi\_bridge\_k}, \texttt{blind}.
          The bridge-width sweep is the paper's centerpiece.
    \item[Axis T --- training.] \texttt{train-free} (mask swapped at
          inference on stock weights) versus \texttt{finetuned} (mask
          present during supervised fine-tuning). Train-free is what
          practitioners can deploy today; fine-tuned is the achievable
          ceiling. The gap between them is a result.
    \item[Axis P --- position layout.] \texttt{sequential} versus
          \texttt{restart}, per Section~\ref{sec:method-positions}. The
          restart layout unlocks the caching co-headline.
    \item[Axis N --- document count.] $\{1, 3, 5, 9, 13, 25\}$. $N=1$ is
          the correctness gate: any DDI variant must be numerically
          identical to \texttt{full} at $N=1$, and any divergence is a
          bug.
    \item[Axis L --- document length.] $\{128, 512, 2048\}$ tokens. Small
          $\ell$ probes the per-block launch overhead where the varlen
          kernel is expected to lose to a dense one; long $\ell$ is where
          DDI's compute win is largest.
    \item[Axis K --- kernel backend.] Dense SDPA with an explicit $-\infty$
          mask (reference semantics), FlashAttention-2 varlen, and a
          Triton block-diagonal implementation. Three implementations of
          one mathematical object, so they double as each other's
          correctness tests.
\end{description}

\subsection{Ground rules}
\label{sec:design-ground-rules}

Five discipline rules apply to every experiment
(issue~\href{https://github.com/AlexKay28/ddi/issues/1}{\#1}):
(1)~every result row is written to the append-only store with full
provenance (git SHA, config, seed, dataset hash, model revision);
(2)~a number that is not in the store cannot appear in the paper --- every
table and figure in Sections~\ref{sec:results-quality}--\ref{sec:failure-envelope}
is rendered from \texttt{results/runs.jsonl} by \texttt{ddi.report};
(3)~one mask definition, one implementation, no re-implementations across
files; (4)~a quality claim requires $\geq 3$ seeds, so per-cell variance
is measurable; (5)~$\Delta$ is fixed from the full baseline before any DDI
number is interpreted (Section~\ref{sec:method-delta}).
Ground rule~(2) is the mechanism by which this preprint's numbers are not
hand-typed: an empty store surfaces as an explicit TODO stub in the
compiled PDF, and no other path exists to place a numerical claim in the
document.
